% Generated draft: requires mathematical and source review.
\documentclass[11pt]{article}
\usepackage[margin=25mm]{geometry}
\usepackage{fontspec}
\setmainfont{Latin Modern Roman}
\usepackage{amsmath,amssymb,mathtools,unicode-math}
\setmathfont{Latin Modern Math}
\usepackage{enumitem,needspace,xurl,hyperref}
\setlength{\parindent}{0pt}
\setlength{\parskip}{5pt}
\setlength{\emergencystretch}{3em}
\newcommand{\N}{\mathbb{N}}
\newcommand{\Z}{\mathbb{Z}}
\newcommand{\Q}{\mathbb{Q}}
\newcommand{\R}{\mathbb{R}}
\newcommand{\C}{\mathbb{C}}
\newcommand{\F}{\mathbb{F}}
\newcommand{\A}{\mathbb{A}}
\newcommand{\PP}{\mathbb{P}}
\newcommand{\E}{\mathbb{E}}
\newcommand{\eps}{\varepsilon}
\newcommand{\dd}{\,\mathrm{d}}
\begin{document}
\section*{UnsolvedMath ID: 800471}
\textbf{Source problem code:} catalogue entry
\subsection{Definitions and mathematical statement}
Let \(N\) denote the ambient space intended by the source, and let \(\mathcal B(N)\) denote its intended class of embedded billiard tables. These are placeholders: the supplied record does not specify \(N\), its geometric structure, or the admissibility conditions for billiard tables. For \(\Omega\in\mathcal B(N)\), retain the source notation \(\Omega\text{-}\mathrm{sys}\) for its unspecified numerical invariant; its definition is absent from the supplied record.

Subject to those unresolved definitions, the question is whether
\[
\exists (\Omega_i)_{i\geq 1}\in\mathcal B(N)^{\{1,2,\ldots\}}
\quad\text{such that}\quad
\Omega_i\text{-}\mathrm{sys}\longrightarrow+\infty
\quad\text{as }i\to\infty.
\]
Here divergence to \(+\infty\) means that, for every real \(M>0\), there exists an integer \(i_0\geq1\) such that \(\Omega_i\text{-}\mathrm{sys}>M\) for every integer \(i\geq i_0\), provided the intended invariant admits these comparisons. No nesting condition on the tables is stated.

The source further says that a positive answer would, to the best of the author's knowledge, give the first example of a disk cotangent bundle with infinite Hofer--Zehnder capacity. This is a supplied, author-qualified consequence and novelty assertion; the record does not identify the relevant bundle or provide the hypotheses or argument linking the question to that consequence.
\subsection{Short English statement}
Does the intended ambient space \(N\) contain a sequence of embedded billiard tables whose source-designated invariant \(\Omega_i\text{-}\mathrm{sys}\) tends to infinity? The ambient geometry, admissible tables, and invariant are unspecified, so this remains a conditional formulation requiring clarification. The registry labels the problem open; that label is supplied metadata, not an independently verified mathematical assessment.
\subsection{Sources}
\begin{itemize}
\item Local UnsolvedMath registry, permanent ID \(800471\), source problem code \texttt{catalogue entry}; supplied complete source record, upstream revision \texttt{d371ee917ac9647fbb06cd6fc529deec6f3dda82}. The record lists Albach, Barraud et al. as proposers, gives the year 2026, and labels the problem open.
\item Albach, Barraud et al., \emph{Open Problems in Billiards and Quantitative Symplectic Geometry} (2026), source item ``Johanna Bimmermann --- Problem 2.'' Supplied reference: \url{https://arxiv.org/abs/2602.12896}. Contributor: Johanna Bimmermann.
\item The supplied provenance records access on 2026-10-16 and extraction by \texttt{arxiv\_latex}. This is a supplied access record, not a claim of independent source consultation or verification. The assertion about a first disk cotangent bundle example with infinite Hofer--Zehnder capacity retains the source author's knowledge qualification.
\end{itemize}
\end{document}
